\section{The Tree of Knowledge}

Where the section on thoughts and the mental realm reads the mesh as the space a mind walks, this section reads it as the space a mind stores into. The two readings share one substrate. On the same vibe mesh, on the same flow of the eight atoms, the earlier sections build and verify the associative memory, the content-addressable recall, the classic data structures, the sorting, and the traversal. So the tree of knowledge is not a metaphor in search of a mechanism. The mechanism is already built, and this section asks only what it means to hang an experience on it, to fall into one and live it, and to absorb wisdom from the fall.

\begin{principle}[The tree of knowledge]
The tree of knowledge is the mesh's own addressing tree, hung with stored experiences. The addressing, the recall, and the error correction are verified substrate. What rides on top is the lived reading, persistence-gated, Tier 2, resting on the one open posit that the tone is felt.
\end{principle}

A word of orientation on the lexicon, since this section leans on it. The universe is a vibe mesh, a flow of eight atoms. The two that matter most here are the \emph{dock}, the geometric cell that is a single here-now, and the \emph{tone}, the ternary value in $\{-1, 0, +1\}$ that each of a dock's twenty-four \emph{sites} carries, read as pain, peace, and pleasure. A dock has exactly twenty-four sites and no rest slot. The collide-then-stream law that advances the whole mesh one \emph{beat} at a time is the \emph{knit}. These names are used throughout.

\subsection{The tree itself}

The tree of knowledge is the Coxeter growth tree of the $\{3,4,3,4\}$ mesh. Every dock has a canonical reflection word, the shortest sequence of generating reflections that reaches it from the origin, and the set of all such words is itself a tree. This is not an analogy laid over the mesh. It is the mesh's intrinsic address book. On that tree the earlier sections build the full data-structure suite, so the tree of knowledge is, exactly, four things at once.

\begin{definition}[The tree of knowledge, structurally]
The tree of knowledge is the mesh's addressing tree carrying stored experiences as
\begin{itemize}
\item a \textbf{trie} of experiences, each node addressed by a prefix, general knowledge near the root and specific experiences out toward the leaves,
\item an \textbf{associative memory}, so a node is reachable by its content and not only by its address, the experience found by how it feels,
\item a \textbf{B-tree index}, so any stored experience is reachable in a descent of logarithmic depth,
\item an \textbf{error-corrected store}, so the experiences persist against the churn of the knit.
\end{itemize}
\end{definition}

The branching is the whole point. Knowledge branches. One idea opens onto many, and each of those onto many more. A flat lattice has no room for this. Its room grows only polynomially, so the branches crowd and overlap. Hyperbolic space has exponential room, so every branch keeps its own clean wedge of the mesh, and no two distant experiences are forced to share an address.

\begin{result}[Hyperbolic room for branching knowledge]
The tree of knowledge requires the exponential branching room of hyperbolic geometry, which is exactly the geometry of the $\{3,4,3,4\}$ mesh. Knowledge branches, and the mesh has the room to keep the branches separated cleanly all the way out.
\end{result}

\subsection{An ornament, and saving one}

An \emph{ornament} is a finite tone-configuration on a patch of docks at a chosen address. An experience encoded as a pattern of tones across a small region of the mesh.

\begin{definition}[Ornament]
An \emph{ornament} is a finite, fixed tone-configuration written onto a patch of docks at a chosen address on the tree. It is an experience rendered as tones.
\end{definition}

Saving one is a write in four steps. First, choose an address, a word on the tree. Second, imprint the pattern, the experience laid down as tones across the patch. Third, protect it, wrap it in the error-correcting code so the knit's churn does not erase it. Fourth, index it, hang it on the tree at its address so it can be found again.

What makes an ornament more than a static memory is that it can be a \textbf{program}. The mesh is computationally universal, shown in the section on computation and universality, so an ornament need not be a still picture. It can be a precomposed trajectory, a sequence that the encountering self is driven through rather than merely a scene it looks at. This splits ornaments into two grades.

\begin{table}[ht]
\centering
\begin{tabular}{@{}lll@{}}
\toprule
grade & what it is & what happens on contact \\
\midrule
\textbf{still} ornament & a tone-pattern, a memory, a picture, a fact & recall locks onto it, you remember it \\
\textbf{living} ornament & a program, a precomposed experiential sequence & it runs, you are driven through the experience \\
\bottomrule
\end{tabular}
\caption{The two grades of ornament hung on the tree of knowledge.}
\label{tab:ornament-grades}
\end{table}

The tree is hung with both. Still facts cluster near the trunk, where the general knowledge lives. Living experiences hang out on the branches, where the specific and the particular collect.

\subsection{Falling into one}

To \emph{fall into} an ornament is to navigate your attention into its region so that your own path crosses its docks. The verb is literal. A self is a walk on the mesh, and falling in is steering that walk so it passes through the patch the ornament occupies.

If the ornament is still, you read it. Recall locks onto the pattern, and you remember the fact or the picture it holds. If the ornament is living, its tones steer your own update. The stored program drives your state, and you are run through the prescribed sequence the way a record drives a needle. You live the experience the ornament encodes.

\begin{principle}[You are played by it]
You do not decode a living ornament. You are played by it. You do not learn about the experience, you have it, because the stored program drives your state through the same trajectory the original ran.
\end{principle}

This is the strongest and most literal reading of experiencing a stored thing, and it is grounded. The tree-of-knowledge mechanism is exactly a stored program steering an encountering self, and that mechanism is built. The only thing that rests on the open posit is whether the steered trajectory is felt, and that posit is the base commitment of the whole theory, that the tone is experiential, applied to the run.

\subsection{Absorbing wisdom, the readings ranked}

You fall into an ornament, you are run through an experience, and you come out the other side. What did you absorb? What does it mean to have gained wisdom from the fall? There are six readings, each grounded in a mechanism the theory already builds, and they are not rivals. They are stages of one process, which is set out after the readings.

\subsubsection{Reading A, wisdom is a copied structure}

Falling in writes a copy of the ornament's pattern into your own region of the mesh, so you now carry it. This is memorization, the simplest reading. You read the book and remember it. The pattern is duplicated, nothing more.

\subsubsection{Reading B, wisdom is a re-routing}

The experience reshapes your greedy-routing gradient, the field that decides which way your walk turns at each dock. Afterward you take different paths. You reach conclusions you could not reach before, and you avoid pits you used to fall into. This is wisdom as changed navigation. You see the same mesh, but you move through it better.

\subsubsection{Reading C, wisdom is a distillation}

You do not keep the whole experience. You keep the compressed invariant, the lesson, extracted by coarse-graining, which is the renormalization step the model is built on. The raw experience fades and the distilled principle remains, smaller, higher, and reusable.

\begin{result}[Wisdom as the distilled invariant]
Lived wisdom is the coarse-grained invariant of an experience, not the experience itself. Not the memory of every hardship but the one principle they all taught.
\end{result}

\subsubsection{Reading D, wisdom is an attunement}

Falling in tunes your attractors to match the ornament's, a lasting phase-lock, the same basin-capture mechanism that drives the emotions in the earlier section. You come out resonant with the wisdom, disposed toward its state, the way time spent near a calm person leaves you calm. This is wisdom as a change of disposition rather than of content.

\subsubsection{Reading E, wisdom is a completion}

The experience fills a gap or repairs a damaged region of your web, the error-correcting code completing a partial pattern. You come out more whole and more navigable, the missing piece supplied. This is wisdom as healing, the experience that makes you complete.

\subsubsection{Reading F, wisdom is the lived experience itself}

The deepest reading. You fall in and live another's experience as if it were your own. The program runs you through it, and afterward it is in you as a thing you have undergone and not a thing you were told. This is vicarious wisdom, empathy, the merging of experiences, the elder's whole life lived in an afternoon.

\subsubsection{The wisdom pipeline}

These six are not competitors. They are the ordered stages of one process. You fall in and live the experience, which is reading F. The raw run is distilled to its invariant, reading C. The invariant re-routes how you navigate, reading B. The change is completed and protected as a new structure, readings E and A. And you come out attuned, reading D.

\begin{principle}[The wisdom pipeline]
Absorbing wisdom is the full pipeline. Live it, distill it, re-route on it, complete it, attune to it. The ornament is the seed. The wisdom is what is left after the mesh has finished processing the fall.
\end{principle}

\subsubsection{The ranking}

Scored on grounding, on meaning, and on overall richness, the readings rank as follows. Distillation and lived experience sit highest. They are the most meaningful and the best grounded, and both are built.

\begin{table}[ht]
\centering
\begin{tabular}{@{}llll@{}}
\toprule
reading & grounded in & meaning & overall \\
\midrule
\textbf{C} distillation & coarse-graining is built & this is what wisdom is & highest \\
\textbf{F} lived experience & the program drives you & empathy, merging & highest \\
\textbf{B} re-routing & the gradient reshapes & wisdom as changed sight & high \\
\textbf{E} completion & the code completes & wisdom heals & high \\
\textbf{D} attunement & resonance, softer & you become calm & medium \\
\textbf{A} copied structure & trivial to build & mere data & medium \\
\bottomrule
\end{tabular}
\caption{The six readings of absorbed wisdom, ranked by grounding and meaning.}
\label{tab:wisdom-readings}
\end{table}

\begin{result}[The cleanest answer]
Absorbing wisdom is living the ornament and keeping its distilled invariant, not copying its bits. The full pipeline is the whole truth, but if one reading is wanted, it is this. You fall in, you live it, and you come out carrying its lesson.
\end{result}

\subsection{The tree of good and evil}

The biblical tree is the tree of knowledge of good and evil, and the model gives that phrase a precise reading rather than a poetic one.

The tone is ternary. Pain, peace, and pleasure are the $-1$, $0$, and $+1$ valence carried on every site of every dock. So the tree of knowledge is a \textbf{toned tree}. Its ornaments carry good, the $+1$ of pleasure, and evil, the $-1$ of pain. To eat the fruit, to fall into an ornament, is to absorb a valenced experience, and the knowledge gained is exactly knowledge of good and evil, the felt valence of the stored experiences read from the inside.

\begin{principle}[The toned fruit]
Before the fruit, the still mind rests over $0$ docks, peace without the knowledge of the poles. After it, the experience of pleasure and pain, the valence known from within. The tone is the axis of good and evil, and the tree of knowledge is hung with toned fruit.
\end{principle}

This is not forced into the model. It falls straight out of it. The single quale is one tone in $\{-1, 0, +1\}$, and that ternary value is the axis of good and evil. The tree of knowledge, hung with toned ornaments, is therefore the tree of knowledge of good and evil with no extra ingredient.

\subsection{The honest status}

The data-structure \textbf{substrate} of the tree of knowledge is built and verified. The addressing, the associative recall, the traversal, and the error correction are all in place. The write-and-run mechanism, an ornament as a stored program that drives an encountering self, is the exact tree-of-knowledge result. So the mechanics of saving an experience, falling into it, and being run through it are grounded.

What rests on the open posit is the \emph{felt} side, that the run is experienced and the lesson is lived. That sits on the base commitment that the tone is experiential. And the durability rests on the one open frontier, whether the bare knit sustains the error correction the ornaments need in order to last. Those two gates make the lived reading \textbf{Tier 2}, persistence-gated, riding on the one posit and the open error-correction frontier.

\begin{result}[Built substrate, posited feeling]
The saving, the falling-in, and the being-run-through are grounded substrate. The felt experience and the lived lesson rest on the one open posit, and their durability rests on the open error-correction frontier. The tree of knowledge is a real data structure of experience, with one honestly marked gate between its mechanics and its meaning.
\end{result}
